\pagenumbering{gobble}

\makecover

\makeabstractpage
{本文围绕黄金、白银、铂金和钯金四类贵金属的期货与现货市场，构建2012年1月至2026年3月的共同交易日价格面板，并计算按100倍缩放的日对数收益率。在此基础上，首先从价格走势、收益率波动、期现相关性和收益率分布特征刻画贵金属期现市场的基本事实；随后使用ICEEMDAN方法对八条收益率序列进行自适应模态分解，并依据IMF的运行长度数采用三成分高斯混合模型将其重构为短期、中期和长期分量。实验结果表明：四类贵金属期货与现货收益率均呈显著正相关，其中铂金和钯金的期现相关性更高；所有收益率序列均表现出负偏、尖峰厚尾和显著平稳性；ICEEMDAN分解能够以接近机器精度的误差重构原序列，GMM-RLN重构进一步揭示了短期高频成分是收益率波动方差的主要来源。上述结果为后续研究贵金属期现市场的多尺度风险联动和跨市场溢出机制提供了数据基础与分解框架。}
{贵金属；期货现货市场；日对数收益率；ICEEMDAN；GMM-RLN；多尺度重构}

\makemaincontents
\makefigtablelist

\startmainmatter

\section{问题描述与研究思路}
贵金属兼具商品属性、金融属性和避险属性，其期货市场承担价格发现与风险管理功能，现货市场则更直接反映即期供需与投资需求变化。近年来，全球通胀周期、货币政策转换、地缘冲突和产业链调整共同推高了贵金属价格波动，使得期货与现货价格之间的同步性、差异性以及收益率风险结构具有较强研究价值。与单一尺度的波动刻画相比，多尺度分解能够将高频冲击、中期调整和长期趋势分离，从而更清晰地识别风险来源。

本文以黄金、白银、铂金和钯金四种贵金属为研究对象，分别考察其期货与现货市场的日度价格和收益率序列。研究思路包括三个步骤：第一，基于原始价格文件构建共同交易日价格面板，并计算日对数收益率；第二，对期货与现货收益率的统计特征进行描述，检验其相关性、非正态性和平稳性；第三，利用ICEEMDAN对各收益率序列进行模态分解，再用GMM-RLN方法将IMF重构为短期、中期和长期分量，为后续风险溢出建模提供多尺度输入。

\section{数据构建与指标说明}
设$P_{i,t}$为第$i$个贵金属市场在交易日$t$的收盘价格，本文按参考文献中常用处理方式计算日对数收益率，并乘以100以便解释：
\begin{equation}
  r_{i,t}=100\left[\ln(P_{i,t})-\ln(P_{i,t-1})\right].
\end{equation}
其中，$i$覆盖黄金、白银、铂金和钯金的期货与现货序列，共8条收益率序列。数据清洗时先分别解析不同来源价格字段，统一日期格式和数值格式，再取所有期货与现货序列的共同非缺失交易日。最终价格面板覆盖2012年1月3日至2026年3月31日，共3101个共同交易日；收益率样本覆盖2012年1月4日至2026年3月31日，共3100个观测值。

\begin{table}[H]
\caption{贵金属期货与现货价格数据来源}
\label{tab:data-source}
\begin{normaltable}
\resizebox{\textwidth}{!}{%
\begin{tabular}{ccccc}
\toprule
品种 & 期货价格文件 & 现货价格文件 & 原始覆盖区间 & 期货/现货原始观测数 \\
\midrule
黄金 & Gold Futures Historical Data.csv & xauusd\_d.csv & 2012-01-03--2026-03-31 & 3257/3675 \\
白银 & Silver Futures Historical Data.csv & xagusd\_d.csv & 2012-01-03--2026-03-31 & 3164/3675 \\
铂金 & Platinum Futures Historical Data.csv & xptusd\_d.csv & 2012-01-02--2026-03-31 & 4161/3675 \\
钯金 & Palladium Futures Historical Data.csv & xpdusd\_d.csv & 2012-01-02--2026-03-31 & 3944/3675 \\
\bottomrule
\end{tabular}}
\end{normaltable}
\end{table}

\section{模型建立、求解与检验}
对任一收益率序列$r_{i,t}$，ICEEMDAN将其分解为若干个由高频到低频排列的IMF分量与残差项：
\begin{equation}
  r_{i,t}=\sum_{k=1}^{K_i} c_{i,k,t}+\varepsilon_{i,t},
\end{equation}
其中，$c_{i,k,t}$表示第$k$个IMF，$\varepsilon_{i,t}$表示残差项。实验中设置噪声试验次数为100，噪声幅度参数为0.2，最大IMF数为12，并在每条序列上检验IMF与残差和对原序列的重构误差。

为形成可用于后续风险网络分析的多尺度分量，本文进一步计算各IMF的运行长度数（Run Length Number, RLN），并用三成分高斯混合模型进行聚类。按聚类中心对应的RLN大小，将IMF划分为短期成分（STC）、中期成分（MTC）和长期成分（LTC）。于是第$i$条序列在尺度$s$上的重构分量可写为
\begin{equation}
  R_{i,t}^{(s)}=\sum_{k\in \mathcal{C}_{i}^{(s)}}c_{i,k,t},\quad s\in\{\mathrm{STC},\mathrm{MTC},\mathrm{LTC}\},
\end{equation}
其中$\mathcal{C}_{i}^{(s)}$表示被GMM-RLN归入尺度$s$的IMF集合。长期成分同时吸收残差项，使得
\begin{equation}
  r_{i,t}=R_{i,t}^{(\mathrm{STC})}+R_{i,t}^{(\mathrm{MTC})}+R_{i,t}^{(\mathrm{LTC})}.
\end{equation}

\section{数据描述}
\subsection{数据来源}
本文选取黄金、白银、铂金和钯金四类贵金属作为研究对象。四类品种均具有活跃的期货交易和国际化现货报价：黄金和白银更受避险需求、美元流动性和通胀预期影响；铂金和钯金除金融属性外，还与汽车催化剂、工业制造等实际需求密切相关。因此，将四类贵金属纳入同一框架，有助于比较贵金属市场中金融属性与产业属性的差异。

期货价格来自各贵金属历史期货价格文件，现货价格来自相应美元计价现货文件。由于不同市场的节假日和交易记录长度并不完全一致，本文不采用单条序列的最长样本，而是保留八条价格序列同时存在的共同交易日。处理后价格面板没有缺失值和重复日期，且共同样本不包含周末交易日。2017年以后共同交易日数量下降至每年约198--201个，因此对低频长期成分的经济解释保持审慎。

\subsection{统计描述}
图\ref{fig:price}展示了四类贵金属期货与现货价格的演化。总体来看，同一品种的期货和现货价格走势高度一致，说明两个市场在长期价格方向上具有较强同步性。黄金价格在样本后段明显上行，体现其避险和金融属性；白银价格波动弹性更强；铂金和钯金价格则呈现更明显的阶段性变化，反映工业需求和供给扰动对其价格的影响。

\begin{figure}[H]
  \centering
  \includegraphics[width=\textwidth]{fig_precious_metals_prices.png}
  \caption{四类贵金属期货与现货价格走势}
  \label{fig:price}
\end{figure}

图\ref{fig:return}给出了对应的日对数收益率。与价格序列相比，收益率围绕零上下波动，并呈现出显著的波动集聚特征，即大幅波动往往集中出现在相邻时期。白银、铂金和钯金的收益率振幅普遍大于黄金，说明黄金相对更稳健，而白银和铂族金属对市场冲击更敏感。

\begin{figure}[H]
  \centering
  \includegraphics[width=\textwidth]{fig_precious_metals_returns.png}
  \caption{四类贵金属期货与现货日对数收益率走势}
  \label{fig:return}
\end{figure}

表\ref{tab:corr}报告了同一品种期货与现货收益率的相关系数。四类贵金属相关系数均高于0.85，其中铂金和钯金分别达到0.927和0.944，说明铂族金属期现市场在日度收益层面联动更紧密。黄金和白银相关系数略低，但仍表现出显著同步性。

\begin{table}[H]
\caption{贵金属期货与现货收益率相关系数}
\label{tab:corr}
\begin{normaltable}
\begin{tabular}{ccc}
\toprule
品种 & 期现收益率相关系数 & 重叠观测数 \\
\midrule
黄金 & 0.854 & 3100 \\
白银 & 0.868 & 3100 \\
铂金 & 0.927 & 3100 \\
钯金 & 0.944 & 3100 \\
\bottomrule
\end{tabular}
\end{normaltable}
\end{table}

表\ref{tab:descriptive}进一步给出收益率描述统计。所有序列均值接近零，符合日度金融收益率的一般特征；标准差方面，黄金最小，钯金和白银较大，表明后两者波动风险更强。各序列偏度均为负，说明样本期内负向极端冲击更突出；峰度均显著大于3，Jarque-Bera统计量对应的$p$值均近似为0，表明收益率分布明显偏离正态分布并具有尖峰厚尾特征。ADF检验统计量均显著为负，且$p$值小于1\%，说明收益率序列平稳，可以用于后续分解和重构分析。

\begin{table}[H]
\caption{贵金属期货与现货收益率描述统计}
\label{tab:descriptive}
\begin{normaltable}
\resizebox{\textwidth}{!}{%
\begin{tabular}{crrrrrrr}
\toprule
序列 & 样本量 & 均值 & 标准差 & 偏度 & 峰度 & JB统计量 & ADF统计量 \\
\midrule
黄金期货 & 3100 & 0.035 & 1.128 & -0.755 & 12.188 & 11158.9 & -57.18 \\
黄金现货 & 3100 & 0.034 & 1.085 & -0.600 & 10.182 & 6823.3 & -40.70 \\
白银期货 & 3100 & 0.030 & 2.222 & -1.739 & 36.856 & 149119.4 & -14.56 \\
白银现货 & 3100 & 0.030 & 2.093 & -1.943 & 33.988 & 125569.6 & -40.56 \\
铂金期货 & 3100 & 0.010 & 1.848 & -0.663 & 13.198 & 13611.5 & -24.53 \\
铂金现货 & 3100 & 0.010 & 1.737 & -0.800 & 13.743 & 15184.4 & -14.44 \\
钯金期货 & 3100 & 0.026 & 2.446 & -0.687 & 12.593 & 12085.9 & -25.08 \\
钯金现货 & 3100 & 0.026 & 2.306 & -0.768 & 13.192 & 13673.5 & -23.22 \\
\bottomrule
\end{tabular}}
\end{normaltable}
\end{table}

\section{实证分析}
\subsection{贵金属收益率序列分解}
参考多尺度风险研究的处理思路，本文首先对八条贵金属期货与现货收益率序列进行ICEEMDAN分解。图\ref{fig:imf}以黄金期货为例展示IMF时间曲线。可以看到，IMF1的振荡频率最高，主要捕捉日度噪声、突发消息和短期交易冲击；随着IMF阶数上升，分量波动逐渐平滑，平均周期逐步拉长，较高阶IMF与残差项更多反映低频调整和长期趋势。

\begin{figure}[H]
  \centering
  \includegraphics[width=\textwidth]{fig_gold_futures_imfs.png}
  \caption{黄金期货收益率的ICEEMDAN IMF时间曲线}
  \label{fig:imf}
\end{figure}

从全部序列看，除钯金期货分解得到9个IMF外，其余七条序列均得到10个IMF。各序列IMF平均周期均从约2.7个交易日开始，随后逐步扩展至1550或3100个交易日，说明ICEEMDAN在样本内形成了由高频到低频的有序分解。所有序列的IMF与残差重构误差均处于$10^{-18}$量级，表明分解在数值上能够精确还原原始收益率。

\begin{table}[H]
\caption{ICEEMDAN分解数量与重构误差}
\label{tab:iceemdan-error}
\begin{normaltable}
\begin{tabular}{cccc}
\toprule
序列 & IMF数量 & IMF重构最大误差 & GMM-RLN重构最大误差 \\
\midrule
黄金期货 & 10 & $6.94\times10^{-18}$ & $1.78\times10^{-15}$ \\
黄金现货 & 10 & $6.94\times10^{-18}$ & $3.55\times10^{-15}$ \\
白银期货 & 10 & $6.94\times10^{-18}$ & $3.55\times10^{-15}$ \\
白银现货 & 10 & $8.67\times10^{-19}$ & $7.11\times10^{-15}$ \\
铂金期货 & 10 & $3.47\times10^{-18}$ & $2.66\times10^{-15}$ \\
铂金现货 & 10 & $3.47\times10^{-18}$ & $1.78\times10^{-15}$ \\
钯金期货 & 9 & $1.73\times10^{-18}$ & $5.33\times10^{-15}$ \\
钯金现货 & 10 & $6.94\times10^{-18}$ & $3.55\times10^{-15}$ \\
\bottomrule
\end{tabular}
\end{normaltable}
\end{table}

\subsection{贵金属收益率序列重构}
在IMF分解基础上，本文使用GMM-RLN将IMF聚类为短期、中期和长期三类尺度。表\ref{tab:gmm-scale}显示，各序列通常由1个短期IMF、1个中期IMF和7--8个长期IMF构成。这一结果表明，高频与中频波动往往集中在最前两个IMF中，而低频波动由多个较平滑的IMF与残差共同刻画。

\begin{table}[H]
\caption{GMM-RLN多尺度分组及方差贡献}
\label{tab:gmm-scale}
\begin{normaltable}
\resizebox{\textwidth}{!}{%
\begin{tabular}{crrrrrr}
\toprule
序列 & STC数量 & STC方差占比(\%) & MTC数量 & MTC方差占比(\%) & LTC数量 & LTC方差占比(\%) \\
\midrule
黄金期货 & 1 & 58.0 & 1 & 20.9 & 8 & 21.0 \\
黄金现货 & 1 & 58.4 & 1 & 19.3 & 8 & 22.2 \\
白银期货 & 1 & 54.8 & 1 & 16.4 & 8 & 28.8 \\
白银现货 & 1 & 53.7 & 1 & 16.8 & 8 & 29.5 \\
铂金期货 & 1 & 61.0 & 1 & 19.0 & 8 & 19.9 \\
铂金现货 & 1 & 61.1 & 1 & 18.7 & 8 & 20.2 \\
钯金期货 & 1 & 55.9 & 1 & 24.3 & 7 & 19.8 \\
钯金现货 & 1 & 54.3 & 1 & 21.8 & 8 & 24.0 \\
\bottomrule
\end{tabular}}
\end{normaltable}
\end{table}

从方差贡献看，短期成分在所有序列中均占据最大比例，约为53.7\%--61.1\%。这说明贵金属日收益率波动主要受短期市场冲击驱动，尤其是铂金期货和铂金现货的短期方差占比超过61\%。中期成分方差占比约为16.4\%--24.3\%，其中钯金期货和钯金现货的中期贡献较高，提示钯金市场可能更容易受到阶段性产业供需调整影响。长期成分虽然由更多IMF构成，但方差占比相对有限，主要承担低频趋势和残差修正功能。

图\ref{fig:reconstruction}以黄金期货为例展示原始收益率、IMF加总重构曲线、STC+MTC+LTC重构曲线及误差。两条重构曲线与原始序列几乎完全重合，误差曲线围绕零且量级极小，验证了ICEEMDAN分解和GMM-RLN重构的数值一致性。由此得到的短期、中期和长期重构序列可作为后续构建多尺度风险溢出网络的直接输入。

\begin{figure}[H]
  \centering
  \includegraphics[width=\textwidth]{fig_gold_futures_reconstruction.png}
  \caption{黄金期货收益率的多尺度重构曲线}
  \label{fig:reconstruction}
\end{figure}

\subsection{实证结果小结}
综合数据描述与多尺度分解结果，可以得到三点结论。第一，四类贵金属期货与现货收益率具有显著联动关系，但不同品种的联动强度存在差异，铂金和钯金高于黄金与白银。第二，收益率序列普遍呈现负偏和尖峰厚尾，说明极端下跌冲击不可忽视，采用能够适应非线性、非平稳局部结构的分解方法具有必要性。第三，GMM-RLN重构结果显示短期成分是贵金属收益率波动的主要来源，长期成分更多反映低频背景变化。因此，在进一步研究贵金属市场风险溢出时，应区分短期冲击传导和长期趋势联动，避免将不同时间尺度的风险机制混合解释。

\section{结论与建议}
本文基于共同交易日价格面板，完成了贵金属期货与现货收益率构造、统计描述、ICEEMDAN分解和GMM-RLN多尺度重构。实证结果说明，贵金属期现市场具有较强同步性，但收益率分布显著偏离正态且存在明显极端波动；ICEEMDAN能够稳定提取由高频到低频的IMF分量，GMM-RLN方法进一步将其压缩为具有经济解释的短期、中期和长期成分。

后续研究可在本文重构结果基础上继续构建多尺度风险溢出网络，比较STC、MTC和LTC三个尺度下期货与现货市场的风险传导方向、中心节点和时变结构。同时，由于2017年以后共同交易日数量有所下降，长期成分解释应结合更多数据源或稳健性检验进一步验证。

\begin{cnreferences}
  \item Huang N E, Shen Z, Long S R, et al. The empirical mode decomposition and the Hilbert spectrum for nonlinear and non-stationary time series analysis[J]. Proceedings of the Royal Society A, 1998, 454(1971): 903--995.
  \item Wu Z, Huang N E. Ensemble empirical mode decomposition: a noise-assisted data analysis method[J]. Advances in Adaptive Data Analysis, 2009, 1(1): 1--41.
  \item Hamilton J D. Time Series Analysis[M]. Princeton: Princeton University Press, 1994.
\end{cnreferences}

\appendixheading
本文主要数据与结果文件由两个实验notebook生成，核心输出如下：
\begin{itemize}[leftmargin=2em,itemsep=0pt,topsep=0pt]
  \item 01号notebook负责价格面板和收益率统计表构造。
  \item 02号notebook负责ICEEMDAN分解、GMM-RLN重构、诊断表和可视化输出。
  \item 核心CSV包括收益率面板、描述统计、ICEEMDAN诊断表和GMM-RLN重构宽表。
\end{itemize}

\begin{table}[H]
\caption{本文核心输出文件说明}
\label{tab:abbr}
\begin{normaltable}
\resizebox{\textwidth}{!}{%
\begin{tabular}{cc}
\toprule
文件 & 内容说明 \\
\midrule
\texttt{precious\_metals\_price\_panel\_complete.csv} & 四类贵金属期货与现货共同交易日价格面板 \\
\texttt{precious\_metals\_daily\_log\_returns.csv} & 八条收益率序列，按$100\Delta\ln P$计算 \\
\texttt{precious\_metals\_return\_descriptive\_stats.csv} & 收益率描述统计、JB检验和ADF检验结果 \\
\texttt{iceemdan\_diagnostics.csv} & IMF平均周期、方差贡献、RLN聚类和尺度标签 \\
\texttt{gmm\_rln\_components\_wide.csv} & STC、MTC、LTC三类重构成分宽表 \\
\bottomrule
\end{tabular}}
\end{normaltable}
\end{table}

\acknowledgementheading
感谢开源Python科学生态为本文数据处理、统计检验、分解建模和可视化提供工具支持。本文全部表格与图形均基于本地实验结果生成，后续版本可在补充更多市场变量后继续扩展。

\hfill \submitdate
